\uuid{SMi0}
\titre{Récurrence}
\auteur{bodin ; hanani ; mzari}
\organisation{exo7}
\licence{CC BY-NC-SA 4.0}
\datecreate{2020-01-20}
\niveau{L1}
\module{Algèbre}
\chapitre{Logique, ensemble, raisonnement}
\sousChapitre{Absurde et contraposée}
\difficulte{2}
\theme{Raisonnement, récurrence}
\exo7qcmid{lille-1.logique.024}

\contenu{
\begin{question}
\qclassification{100.03,100.04}
\qtags{motcle=récurrence}
Je souhaite montrer par récurrence une certaine assertion $H_n$, pour tout entier $n\ge0$.
Quels sont les débuts valables pour la rédaction de l'étape d'hérédité ?
\begin{answers}
    \bad{Je suppose $H_n$ vraie pour tout $n\ge0$, et je montre que $H_{n+1}$ est vraie.}

    \bad{Je suppose $H_{n-1}$ vraie pour tout $n\ge1$, et je montre que $H_{n}$ est vraie.}

    \good{Je fixe $n\ge0$, je suppose $H_n$ vraie, et je montre que $H_{n+1}$ est vraie.}

    \bad{Je fixe $n\ge0$ et je montre que $H_{n+1}$ est vraie.}
\end{answers}
\begin{explanations}
La récurrence a une rédaction très rigide. Sinon on raconte vite n'importe quoi !
\end{explanations}
\end{question}
}
